\documentclass[10pt,a4paper]{amsart}
\usepackage{amsmath,amssymb,amsthm,graphicx}
\usepackage{verbatim}
\newenvironment{reference}{\comment}{\endcomment}
%\newenvironment{reference}{}{}
\newenvironment{slogan}{\comment}{\endcomment}
\newenvironment{history}{\comment}{\endcomment}

% For commutative diagrams we use Xy-pic
\usepackage[all]{xy}

% We use 2cell for 2-commutative diagrams.
\xyoption{2cell}
% Theorem environments.
%
\theoremstyle{plain}
\newtheorem{theorem}[subsection]{Theorem}
\newtheorem{proposition}[subsection]{Proposition}
\newtheorem{lemma}[subsection]{Lemma}

\theoremstyle{definition}
\newtheorem{definition}[subsection]{Definition}
\newtheorem{example}[subsection]{Example}
\newtheorem{exercise}[subsection]{Exercise}
\newtheorem{situation}[subsection]{Situation}

\theoremstyle{remark}
\newtheorem{remark}[subsection]{Remark}
\newtheorem{remarks}[subsection]{Remarks}

\numberwithin{equation}{subsection}

% Macros
%
\def\lim{\mathop{\mathrm{lim}}\nolimits}
\def\colim{\mathop{\mathrm{colim}}\nolimits}
\def\Spec{\mathop{\mathrm{Spec}}}
\def\Hom{\mathop{\mathrm{Hom}}\nolimits}
\def\Ext{\mathop{\mathrm{Ext}}\nolimits}
\def\SheafHom{\mathop{\mathcal{H}\!\mathit{om}}\nolimits}
\def\SheafExt{\mathop{\mathcal{E}\!\mathit{xt}}\nolimits}
\def\Sch{\mathit{Sch}}
\def\Mor{\mathop{\mathrm{Mor}}\nolimits}
\def\Ob{\mathop{\mathrm{Ob}}\nolimits}
\def\Sh{\mathop{\mathit{Sh}}\nolimits}
\def\NL{\mathop{N\!L}\nolimits}
\def\CH{\mathop{\mathrm{CH}}\nolimits}
\def\proetale{{pro\text{-}\acute{e}tale}}
\def\etale{{\acute{e}tale}}
\def\QCoh{\mathit{QCoh}}
\def\Ker{\mathop{\mathrm{Ker}}}
\def\Im{\mathop{\mathrm{Im}}}
\def\Coker{\mathop{\mathrm{Coker}}}
\def\Coim{\mathop{\mathrm{Coim}}}

% Boxtimes
%
\DeclareMathSymbol{\boxtimes}{\mathbin}{AMSa}{"02}

%
% Macros for moduli stacks/spaces
%
\def\QCohstack{\mathcal{QC}\!\mathit{oh}}
\def\Cohstack{\mathcal{C}\!\mathit{oh}}
\def\Spacesstack{\mathcal{S}\!\mathit{paces}}
\def\Quotfunctor{\mathrm{Quot}}
\def\Hilbfunctor{\mathrm{Hilb}}
\def\Curvesstack{\mathcal{C}\!\mathit{urves}}
\def\Polarizedstack{\mathcal{P}\!\mathit{olarized}}
\def\Complexesstack{\mathcal{C}\!\mathit{omplexes}}
% \Pic is the operator that assigns to X its picard group, usage \Pic(X)
% \Picardstack_{X/B} denotes the Picard stack of X over B
% \Picardfunctor_{X/B} denotes the Picard functor of X over B
\def\Pic{\mathop{\mathrm{Pic}}\nolimits}
\def\Picardstack{\mathcal{P}\!\mathit{ic}}
\def\Picardfunctor{\mathrm{Pic}}
\def\Deformationcategory{\mathcal{D}\!\mathit{ef}}
\usepackage{cleveref}
\crefname{originallemma}{original lemma}{original lemmas}
\Crefname{originallemma}{Original lemma}{Original lemmas}
\makeatletter
\newlabel{algebra-lemma-complete-local-Noetherian-domain-finite-over-regular}{{032D}{source}{}{sourcepointer-1}{}}
\newlabel{algebra-lemma-complete-local-Noetherian-domain-finite-over-regular@cref}{{[originallemma][0][]032D}{[1][0][]source}{}{}{}}
\makeatother
\crefname{originalproposition}{original proposition}{original propositions}
\Crefname{originalproposition}{Original proposition}{Original propositions}
\makeatletter
\newlabel{algebra-proposition-ubiquity-nagata}{{0335}{source}{}{sourcepointer-2}{}}
\newlabel{algebra-proposition-ubiquity-nagata@cref}{{[originalproposition][0][]0335}{[1][0][]source}{}{}{}}
\makeatother
\crefname{originallemma}{original lemma}{original lemmas}
\Crefname{originallemma}{Original lemma}{Original lemmas}
\makeatletter
\newlabel{dualizing-lemma-ubiquity-dualizing}{{0BFR}{source}{}{sourcepointer-3}{}}
\newlabel{dualizing-lemma-ubiquity-dualizing@cref}{{[originallemma][0][]0BFR}{[1][0][]source}{}{}{}}
\makeatother
\crefname{originalreference}{original reference}{original references}
\Crefname{originalreference}{Original reference}{Original references}
\makeatletter
\newlabel{lemma-equivalence-sequence-blowups}{{0BFY}{source}{}{sourcepointer-4}{}}
\newlabel{lemma-equivalence-sequence-blowups@cref}{{[originalreference][0][]0BFY}{[1][0][]source}{}{}{}}
\makeatother
\crefname{originalreference}{original reference}{original references}
\Crefname{originalreference}{Original reference}{Original references}
\makeatletter
\newlabel{lemma-equivalence-sequence-normalized-blowups}{{0BFZ}{source}{}{sourcepointer-5}{}}
\newlabel{lemma-equivalence-sequence-normalized-blowups@cref}{{[originalreference][0][]0BFZ}{[1][0][]source}{}{}{}}
\makeatother
\crefname{originalreference}{original reference}{original references}
\Crefname{originalreference}{Original reference}{Original references}
\makeatletter
\newlabel{lemma-existence-implies-existence-by-normalized-blowing-ups}{{0BGM}{source}{}{sourcepointer-6}{}}
\newlabel{lemma-existence-implies-existence-by-normalized-blowing-ups@cref}{{[originalreference][0][]0BGM}{[1][0][]source}{}{}{}}
\makeatother
\crefname{originalreference}{original reference}{original references}
\Crefname{originalreference}{Original reference}{Original references}
\makeatletter
\newlabel{lemma-go-up-degree-p}{{0B4U}{source}{}{sourcepointer-7}{}}
\newlabel{lemma-go-up-degree-p@cref}{{[originalreference][0][]0B4U}{[1][0][]source}{}{}{}}
\makeatother
\crefname{originalreference}{original reference}{original references}
\Crefname{originalreference}{Original reference}{Original references}
\makeatletter
\newlabel{lemma-go-up-separable}{{0AXL}{source}{}{sourcepointer-8}{}}
\newlabel{lemma-go-up-separable@cref}{{[originalreference][0][]0AXL}{[1][0][]source}{}{}{}}
\makeatother
\crefname{originalreference}{original reference}{original references}
\Crefname{originalreference}{Original reference}{Original references}
\makeatletter
\newlabel{lemma-normalized-blowup-completion}{{0BGA}{source}{}{sourcepointer-9}{}}
\newlabel{lemma-normalized-blowup-completion@cref}{{[originalreference][0][]0BGA}{[1][0][]source}{}{}{}}
\makeatother
\crefname{originalreference}{original reference}{original references}
\Crefname{originalreference}{Original reference}{Original references}
\makeatletter
\newlabel{lemma-port-regularity-to-completion}{{0BG6}{source}{}{sourcepointer-10}{}}
\newlabel{lemma-port-regularity-to-completion@cref}{{[originalreference][0][]0BG6}{[1][0][]source}{}{}{}}
\makeatother
\crefname{originalreference}{original reference}{original references}
\Crefname{originalreference}{Original reference}{Original references}
\makeatletter
\newlabel{lemma-rational-to-gorenstein}{{0BBV}{source}{}{sourcepointer-11}{}}
\newlabel{lemma-rational-to-gorenstein@cref}{{[originalreference][0][]0BBV}{[1][0][]source}{}{}{}}
\makeatother
\crefname{originalreference}{original reference}{original references}
\Crefname{originalreference}{Original reference}{Original references}
\makeatletter
\newlabel{lemma-reduce-to-rational}{{0B4P}{source}{}{sourcepointer-12}{}}
\newlabel{lemma-reduce-to-rational@cref}{{[originalreference][0][]0B4P}{[1][0][]source}{}{}{}}
\makeatother
\crefname{originalreference}{original reference}{original references}
\Crefname{originalreference}{Original reference}{Original references}
\makeatletter
\newlabel{lemma-resolve-rational-double-points}{{0BGE}{source}{}{sourcepointer-13}{}}
\newlabel{lemma-resolve-rational-double-points@cref}{{[originalreference][0][]0BGE}{[1][0][]source}{}{}{}}
\makeatother
\crefname{originallemma}{original lemma}{original lemmas}
\Crefname{originallemma}{Original lemma}{Original lemmas}
\makeatletter
\newlabel{more-algebra-lemma-normal-goes-up}{{0BFK}{source}{}{sourcepointer-14}{}}
\newlabel{more-algebra-lemma-normal-goes-up@cref}{{[originallemma][0][]0BFK}{[1][0][]source}{}{}{}}
\makeatother
\crefname{originalproposition}{original proposition}{original propositions}
\Crefname{originalproposition}{Original proposition}{Original propositions}
\makeatletter
\newlabel{more-algebra-proposition-ubiquity-J-2}{{07PJ}{source}{}{sourcepointer-15}{}}
\newlabel{more-algebra-proposition-ubiquity-J-2@cref}{{[originalproposition][0][]07PJ}{[1][0][]source}{}{}{}}
\makeatother
\crefname{originalproposition}{original proposition}{original propositions}
\Crefname{originalproposition}{Original proposition}{Original propositions}
\makeatletter
\newlabel{more-algebra-proposition-ubiquity-excellent}{{07QW}{source}{}{sourcepointer-16}{}}
\newlabel{more-algebra-proposition-ubiquity-excellent@cref}{{[originalproposition][0][]07QW}{[1][0][]source}{}{}{}}
\makeatother
\title{A normal two dimensional complete Noetherian local ring admits a resolution of singularities}
\author{The Stacks Project Authors}
\date{stacks-project:resolve.tex:local-snapshot-20261012; GNU FDL 1.2 or later}
\begin{document}
\maketitle
\noindent\textit{Editorial note.} All delivered mathematical text is the author's own, copied from the pinned Stacks Project chapter file resolve.tex; no statement or proof step was rewritten, reordered or removed.\par\medskip
\noindent\textit{Editorial note.} The counted proof is the author's complete argument: the excellency, Nagata and dualizing input together with the normality of completions, the induction on the degree of the fraction field extension, the reduction through an intermediate field with normalization of the base change and resolutions of the completed local rings, and the conversion into normalized blowups followed by the resolution of rational double points.\par\medskip
\noindent\textit{Editorial note.} The label lemma-resolve-complete names the counted statement, and the proof environment immediately following it in the source is delivered verbatim as its proof.\par\medskip
\noindent\textit{Editorial note.} The counted statement is carried by the author's own lemma environment in the chapter source.\par\medskip
\noindent\textit{Editorial note.} Pre-open gate: python operations/agent-authoring/tools/precard\_gate.py --source resolve.tex --label lemma-resolve-complete --start 4050 --end 4156 --local-id 28068 returned RESULT=OK (L28068 free, occupancy FREE, public interval clear). Candidate taken from sources/stacks-resolve/candidates-halfA.json (sixth pool, assigned by the Lead on 2026-10-13).\par\medskip
\noindent\textit{Editorial note.} The statement and proof are copied from the pinned chapter file; the cited inputs are kept as external pointers with their canonical Stacks tags.\par\medskip
\noindent\textit{Allowed standard prerequisites (editorial).}\begin{itemize}
\item Excellency properties of complete Noetherian local rings. More on Algebra, Algebra and Dualizing Complexes, the ubiquity statements that a complete Noetherian local ring is J-2, Nagata and excellent and has a dualizing complex, and the lemma that normality ascends along a regular homomorphism. Supplies the excellency properties of the rings encountered, including those essentially of finite type over the given ring, and the normality of completions.
\item Finiteness of a normal domain over a regular complete local ring. Algebra, the lemma producing a regular complete local subring over which a complete Noetherian local domain is finite. Supplies the finite extension A\_0 -> A of regular complete local rings on whose fraction field degree the induction runs.
\item Replacement of a resolution of a completion by normalized blowups of the original. Same chapter, the lemma turning the existence of resolutions over the completions of the local rings into existence over the original, phrased through normalized blowups. Supplies the passage from resolutions of the completed local rings to a sequence of normalized blowups of the surface.
\item Blowups preserving normality with Gorenstein local rings, and the resolution of rational double points. Same chapter, the lemma that the normalized blowups of the sequence combine into a single sequence of blowups at Gorenstein points, and the lemma resolving the singularities that appear there. Supply the intermediate surface with Gorenstein local rings having only rational double point singularities, which are then resolved.
\end{itemize}\medskip
\ifdefined\hypertarget\hypertarget{sourcepointer-1}{}\fi
\noindent\textit{Source reference.} Original lemma 032D: sources/stacks-chow/tags.txt, source lines 4016--4016. Prerequisite pointer invoked at proof line(s) 4075; it names a standard fact of a sibling chapter that the counted argument uses, and every step of the counted proof built on it is retained verbatim in the delivered body.\par\medskip
\ifdefined\hypertarget\hypertarget{sourcepointer-2}{}\fi
\noindent\textit{Source reference.} Original proposition 0335: sources/stacks-chow/tags.txt, source lines 4044--4044. Prerequisite pointer invoked at proof line(s) 4060; it names a standard fact of a sibling chapter that the counted argument uses, and every step of the counted proof built on it is retained verbatim in the delivered body.\par\medskip
\ifdefined\hypertarget\hypertarget{sourcepointer-3}{}\fi
\noindent\textit{Source reference.} Original lemma 0BFR: sources/stacks-chow/tags.txt, source lines 14394--14394. Prerequisite pointer invoked at proof line(s) 4064; it names a standard fact of a sibling chapter that the counted argument uses, and every step of the counted proof built on it is retained verbatim in the delivered body.\par\medskip
\ifdefined\hypertarget\hypertarget{sourcepointer-4}{}\fi
\noindent\textit{Source reference.} Original reference 0BFY: sources/stacks-resolve/pinned-original-source.tex, source lines 1495--1495. Prerequisite pointer invoked at proof line(s) 4143; it names a standard fact of the same chapter that the counted argument uses, and every step of the counted proof built on it is retained verbatim in the delivered body.\par\medskip
\ifdefined\hypertarget\hypertarget{sourcepointer-5}{}\fi
\noindent\textit{Source reference.} Original reference 0BFZ: sources/stacks-resolve/pinned-original-source.tex, source lines 1541--1541. Prerequisite pointer invoked at proof line(s) 4116; it names a standard fact of the same chapter that the counted argument uses, and every step of the counted proof built on it is retained verbatim in the delivered body.\par\medskip
\ifdefined\hypertarget\hypertarget{sourcepointer-6}{}\fi
\noindent\textit{Source reference.} Original reference 0BGM: sources/stacks-resolve/pinned-original-source.tex, source lines 3950--3950. Prerequisite pointer invoked at proof line(s) 4102; it names a standard fact of the same chapter that the counted argument uses, and every step of the counted proof built on it is retained verbatim in the delivered body.\par\medskip
\ifdefined\hypertarget\hypertarget{sourcepointer-7}{}\fi
\noindent\textit{Source reference.} Original reference 0B4U: sources/stacks-resolve/pinned-original-source.tex, source lines 2349--2349. Prerequisite pointer invoked at proof line(s) 4129; it names a standard fact of the same chapter that the counted argument uses, and every step of the counted proof built on it is retained verbatim in the delivered body.\par\medskip
\ifdefined\hypertarget\hypertarget{sourcepointer-8}{}\fi
\noindent\textit{Source reference.} Original reference 0AXL: sources/stacks-resolve/pinned-original-source.tex, source lines 2139--2139. Prerequisite pointer invoked at proof line(s) 4129; it names a standard fact of the same chapter that the counted argument uses, and every step of the counted proof built on it is retained verbatim in the delivered body.\par\medskip
\ifdefined\hypertarget\hypertarget{sourcepointer-9}{}\fi
\noindent\textit{Source reference.} Original reference 0BGA: sources/stacks-resolve/pinned-original-source.tex, source lines 3236--3236. Prerequisite pointer invoked at proof line(s) 4109; it names a standard fact of the same chapter that the counted argument uses, and every step of the counted proof built on it is retained verbatim in the delivered body.\par\medskip
\ifdefined\hypertarget\hypertarget{sourcepointer-10}{}\fi
\noindent\textit{Source reference.} Original reference 0BG6: sources/stacks-resolve/pinned-original-source.tex, source lines 2997--2997. Prerequisite pointer invoked at proof line(s) 4115; it names a standard fact of the same chapter that the counted argument uses, and every step of the counted proof built on it is retained verbatim in the delivered body.\par\medskip
\ifdefined\hypertarget\hypertarget{sourcepointer-11}{}\fi
\noindent\textit{Source reference.} Original reference 0BBV: sources/stacks-resolve/pinned-original-source.tex, source lines 2688--2688. Prerequisite pointer invoked at proof line(s) 4136; it names a standard fact of the same chapter that the counted argument uses, and every step of the counted proof built on it is retained verbatim in the delivered body.\par\medskip
\ifdefined\hypertarget\hypertarget{sourcepointer-12}{}\fi
\noindent\textit{Source reference.} Original reference 0B4P: sources/stacks-resolve/pinned-original-source.tex, source lines 2087--2087. Prerequisite pointer invoked at proof line(s) 4131; it names a standard fact of the same chapter that the counted argument uses, and every step of the counted proof built on it is retained verbatim in the delivered body.\par\medskip
\ifdefined\hypertarget\hypertarget{sourcepointer-13}{}\fi
\noindent\textit{Source reference.} Original reference 0BGE: sources/stacks-resolve/pinned-original-source.tex, source lines 3716--3716. Prerequisite pointer invoked at proof line(s) 4154; it names a standard fact of the same chapter that the counted argument uses, and every step of the counted proof built on it is retained verbatim in the delivered body.\par\medskip
\ifdefined\hypertarget\hypertarget{sourcepointer-14}{}\fi
\noindent\textit{Source reference.} Original lemma 0BFK: sources/stacks-chow/tags.txt, source lines 14388--14388. Prerequisite pointer invoked at proof line(s) 4068; it names a standard fact of a sibling chapter that the counted argument uses, and every step of the counted proof built on it is retained verbatim in the delivered body.\par\medskip
\ifdefined\hypertarget\hypertarget{sourcepointer-15}{}\fi
\noindent\textit{Source reference.} Original proposition 07PJ: sources/stacks-chow/tags.txt, source lines 9799--9799. Prerequisite pointer invoked at proof line(s) 4059; it names a standard fact of a sibling chapter that the counted argument uses, and every step of the counted proof built on it is retained verbatim in the delivered body.\par\medskip
\ifdefined\hypertarget\hypertarget{sourcepointer-16}{}\fi
\noindent\textit{Source reference.} Original proposition 07QW: sources/stacks-chow/tags.txt, source lines 9846--9846. Prerequisite pointer invoked at proof line(s) 4062; it names a standard fact of a sibling chapter that the counted argument uses, and every step of the counted proof built on it is retained verbatim in the delivered body.\par\medskip
\begin{lemma}
\label{lemma-resolve-complete}
Let $(A, \mathfrak m, \kappa)$ be a Noetherian complete local ring.
Assume $A$ is a normal domain of dimension $2$. Then $\Spec(A)$ has a
resolution of singularities.
\end{lemma}

\begin{proof}
A Noetherian complete local ring is J-2
(More on Algebra, Proposition \ref{more-algebra-proposition-ubiquity-J-2}),
Nagata (Algebra, Proposition \ref{algebra-proposition-ubiquity-nagata}),
excellent (More on Algebra, Proposition
\ref{more-algebra-proposition-ubiquity-excellent}),
and has a dualizing complex
(Dualizing Complexes, Lemma \ref{dualizing-lemma-ubiquity-dualizing}).
Moreover, the same is true for any ring essentially of finite type over $A$.
If $B$ is a excellent local normal domain, then the completion
$B^\wedge$ is normal (as $B \to B^\wedge$ is regular and
More on Algebra, Lemma \ref{more-algebra-lemma-normal-goes-up} applies).
In other words, the local rings which we encounter in the rest of the proof
will have the required ``excellency'' properties required of them.

\medskip\noindent
Choose $A_0 \subset A$ with $A_0$ a regular complete local ring
and $A_0 \to A$ finite, see Algebra, Lemma
\ref{algebra-lemma-complete-local-Noetherian-domain-finite-over-regular}.
This induces a finite extension of fraction fields $K/K_0$.
We will argue by induction on $[K : K_0]$. The base case is
when the degree is $1$ in which case $A_0 = A$ and the result is true.

\medskip\noindent
Suppose there is an intermediate field $K_0 \subset L \subset K$,
$K_0 \not = L \not = K$. Let $B \subset A$ be the integral closure
of $A_0$ in $L$. By induction we choose a resolution of singularities
$Y \to \Spec(B)$. Let $X$ be the normalization
of $Y \times_{\Spec(B)} \Spec(A)$. Picture:
$$
\xymatrix{
X \ar[r] \ar[d] & \Spec(A) \ar[d] \\
Y \ar[r] & \Spec(B)
}
$$
Since $A$ is J-2 the regular locus of $X$ is open. Since $X$
is a normal surface we conclude that $X$ has at worst finitely
many singular points $x_1, \ldots, x_n$ which are closed points with
$\dim(\mathcal{O}_{X, x_i}) = 2$.
For each $i$ let $y_i \in Y$ be the image.
Since
$\mathcal{O}_{Y, y_i}^\wedge \to \mathcal{O}_{X, x_i}^\wedge$
is finite of smaller degree than before we conclude by
induction hypothesis that $\mathcal{O}_{X, x_i}^\wedge$
has resolution of singularities. By
Lemma \ref{lemma-existence-implies-existence-by-normalized-blowing-ups}
there is a sequence
$$
Z^\wedge_{i, n_i} \to \ldots \to Z^\wedge_{i, 1} \to
\Spec(\mathcal{O}_{X, x_i}^\wedge)
$$
of normalized blowups with $Z^\wedge_{i, n_i}$ regular.
By Lemma \ref{lemma-normalized-blowup-completion}
there is a corresponding sequence of normalized blowing ups
$$
Z_{i, n_i} \to \ldots \to Z_{i, 1} \to \Spec(\mathcal{O}_{X, x_i})
$$
Then $Z_{i, n_i}$ is a regular scheme by
Lemma \ref{lemma-port-regularity-to-completion}.
By Lemma \ref{lemma-equivalence-sequence-normalized-blowups}
we can fit these normalized blowing ups
into a corresponding sequence
$$
Z_n \to Z_{n - 1} \to \ldots \to Z_1 \to X
$$
and of course $Z_n$ is regular too (look at the local rings).
This proves the induction step.

\medskip\noindent
Assume there is no intermediate field $K_0 \subset L \subset K$
with $K_0 \not = L \not = K$. Then either $K/K_0$ is separable
or the characteristic to $K$ is $p$ and $[K : K_0] = p$.
Then either Lemma \ref{lemma-go-up-separable} or \ref{lemma-go-up-degree-p}
implies that reduction to rational singularities is possible.
By Lemma \ref{lemma-reduce-to-rational} we conclude that there exists a
normal modification $X \to \Spec(A)$ such that for
every singular point $x$ of $X$ the local ring $\mathcal{O}_{X, x}$
defines a rational singularity. Since $A$ is J-2 we find that $X$ has
finitely many singular points $x_1, \ldots, x_n$.
By Lemma \ref{lemma-rational-to-gorenstein}
there exists a finite sequence of blowups in singular closed points
$$
X_{i, n_i} \to X_{i, n_i - 1} \to \ldots \to \Spec(\mathcal{O}_{X, x_i})
$$
such that $X_{i, n_i}$ is Gorenstein, i.e., has an
invertible dualizing module. By (the essentially trivial)
Lemma \ref{lemma-equivalence-sequence-blowups}
with $n = \sum n_a$ these sequences correspond to a sequence of
blowups
$$
X_n \to X_{n - 1} \to \ldots \to X
$$
such that $X_n$ is normal and the local rings of $X_n$ are Gorenstein.
Again $X_n$ has a finite number of singular points
$x'_1, \ldots, x'_s$, but this time the singularities are
rational double points, more precisely, the local rings
$\mathcal{O}_{X_n, x'_i}$ are as in
Lemma \ref{lemma-resolve-rational-double-points}.
Arguing exactly as above we conclude that the lemma is true.
\end{proof}

\end{document}
